\scanpage{70}
% archive sheet number: 69
On vérifie que $f$ est \unsure{invariante|invariant} par $\underline{a}$, et on a
\[
\deleted{f(o)=\frac{-\alpha^2-\alpha\beta+2(\alpha+\beta)}{4(\alpha+\beta)^2}}
\qquad\boxed{f(o)=\frac{2-\alpha}{4(\alpha+\beta)}}.
\]
C'est nul ssi $\alpha=2$ (i.e., \inserted{en car. $p>0$,} ssi le car. est \unsure{égal au} \deleted{\unsure{nb}} \unsure{ordre} \deleted{\unsure{entier}} des \unsure{sommets}) : dans ce cas bien sûr, la qu. \unsure{d'origine} circonscrite est \uline{singulière}\ldots)

\diagram{\(\begin{array}{ll|l}
\begin{gathered}
\text{\uline{Cas tétraèdre}}\quad\alpha=\beta=-1\\
\delta'=2\quad\delta=-3
\end{gathered}
&\begin{gathered}
\text{centre : à l'}\infty\text{ ssi car. }2\\
f_0\text{ singulier ssi car }2\\
f\text{ singulier ssi car }\deleted{\unsure{2}}3
\end{gathered}
&\begin{gathered}
Q_{\unsure{E|\overline E}}\text{ birégulière ssi}\\
\text{car. }\ne2,3\\
\text{(alors le centre est}\\
\text{à dist. finie)}
\end{gathered}
\end{array}\)}

\diagram{\(\begin{array}{ll|l}
\begin{gathered}
\deleted{\text{\uline{\unsure{Cas octaèdre}}}}\ \inserted{\text{cube}}\quad\alpha=-1,\beta=0\\
\delta'=+4,\quad\delta=-12
\end{gathered}
&\begin{gathered}
\text{centre : à l'}\infty\text{ ssi car. }2\\
f_0\text{ singulier ssi car }2\\
f\text{ singulier ssi car }2,3
\end{gathered}
&\begin{gathered}
Q_E\text{ birég ssi}\\
\text{car}\ne2,3\\
\text{(alors centre à dist.}\\
\text{finie)}
\end{gathered}
\end{array}\)}

\diagram{\(\begin{array}{ll|l}
\begin{gathered}
\deleted{\text{\uline{cube}}}\ \inserted{\text{octaèdre}}\quad\alpha=0,\beta=-1\\
\delta'=2,\quad\delta=-4
\end{gathered}
&\begin{gathered}
\text{centre à l'}\infty\text{ ssi car. }2\\
f_0\text{ sing. ssi car. }2\\
f\text{ sing. ssi car. }2
\end{gathered}
&\begin{gathered}
Q_E\text{ birég ssi}\\
\text{car. }\ne2\\
\text{(alors centre à}\\
\text{dist. finie)}
\end{gathered}
\end{array}\)}
% Curved arrow beside the two corrected case names, directed upward.
\diagram{\begin{tikzcd}[ampersand replacement=\&]
{}\\{}\arrow[u,bend left=45]
\end{tikzcd}}

\diagram{\(\begin{array}{ll|l}
\begin{gathered}
\text{\uline{icosaèdre}}\quad\left\{\begin{aligned}
&\alpha\text{ satisfaisant }\alpha^2+\alpha-1=0\\
&\beta=-1
\end{aligned}\right.\\
\delta'=1\deleted{\unsure{\alpha}},\quad\delta=-\alpha-3\\
N(\delta)=5
\end{gathered}
&\begin{gathered}
\text{centre à l'}\infty\text{ ssi car. }2\\
f_0\text{ tjrs régulier}\\
f\text{ sing. ssi car. }5
\end{gathered}
&\begin{gathered}
Q_E\text{ birég ssi}\\
\text{car. }\ne5\\
\text{(le centre peut}\\
\text{être à l'}\infty\text{ !)}
\end{gathered}
\end{array}\)}

\diagram{\(\begin{array}{ll|l}
\begin{gathered}
\text{\uline{dodécaèdre}}\quad\alpha=-1,\ \beta\text{ satisfaisant }\deleted{\unsure{\beta^2+\beta=0}}\\
\beta^2+\beta-1=0\\
\deleted{\delta'=\unsure{\beta^2+\beta}}\quad\deleted{N(\delta')=\unsure{5|9}\,!}\\
\deleted{\unsure{\delta}=}\quad\delta'=\beta+2\quad N(\delta')=1\\
\delta=-3(\unsure{\beta+1|\beta-1|\beta+2})^2\quad N(\delta)=9\\
=-9\beta-15
\end{gathered}
&\begin{gathered}
\text{centre à l'}\infty\text{ ssi car }2\\
f_0\text{ tjrs régulier}\\
f\text{ sing. ssi car }3
\end{gathered}
&\begin{gathered}
Q_E\text{ bir}\unsure{\text{ég}}\text{ ssi}\\
\text{car. }\ne3\\
\text{(le centre peut}\\
\text{être à l'}\infty\text{)}
\end{gathered}
\end{array}\)}

\diagram{\(\begin{array}{ll|l}
\begin{gathered}
\text{\uline{\unsure{pentagramme|pentagone}}}\quad\alpha,\beta\text{ solutions de}\\
\lambda^2+\lambda-1=0,\quad\alpha\beta=-1\\
\beta=-1/\alpha=-1-\alpha\\
\delta'=\deleted{\unsure{\beta}}\alpha+2\quad N(\delta')=1\\
\delta=4\alpha-3\quad N(\delta)=5
\end{gathered}
&\begin{gathered}
\text{centre : à l'}\infty\text{ ssi car }2\\
f_0\text{ tjrs régulier}\\
f\text{ singulier ssi car }5
\end{gathered}
&\begin{gathered}
Q_E\text{ birég ssi}\\
\text{car. }\ne5\\
\text{(le centre peut}\\
\text{être à l'}\infty\text{)}
\end{gathered}
\end{array}\)}

\margin{\unsure{?}}
Il y a \deleted{\unsure{trois}} \unsure{enfin|encore} un pol. rég. \inserted{fini} \uline{birégulier} \deleted{(\unsure{en dehors de ceux-là})} mais son dual ne l'est pas \unsure{c'est le cas|c'est} en \unsure{car.} $\ne0$ en car. $3\deleted{\unsure{,5,6}}$ (cube), $3$ (ico), $5$ (dodéc).
